\begin{lemma}[Strict growth in a finite vertex set reaches a terminal state]
Let $V$ be finite, and let $\mathcal S$ and $\mathcal T$ be arbitrary
predicates on subsets of $V$.  Suppose that $\mathcal S(X_0)$ for some
$X_0\subseteq V$, and that for every $X\subseteq V$ satisfying
$\mathcal S(X)$ but not $\mathcal T(X)$ there is a set $Y\subseteq V$
such that $\mathcal S(Y)$ and $X\subsetneq Y$.  Then there is a set
$X\subseteq V$ satisfying both $\mathcal S(X)$ and $\mathcal T(X)$.
\end{lemma}
\begin{proof}
For a subset $X\subseteq V$, define its deficit by
$d(X)=|V|-|X|$.  We prove by well-founded induction on the natural number
$d(X)$ that every $X$ satisfying $\mathcal S(X)$ has a terminal successor
$Z$ satisfying both $\mathcal S(Z)$ and $\mathcal T(Z)$.

Fix $X$ with $\mathcal S(X)$.  If $\mathcal T(X)$ holds, take $Z=X$.
Otherwise the growth hypothesis supplies $Y$ with $\mathcal S(Y)$ and
$X\subsetneq Y$.  Strict inclusion of finite sets gives $|X|<|Y|$, while
$Y\subseteq V$ gives $|Y|\leq |V|$.  Therefore
\[
              d(Y)=|V|-|Y|<|V|-|X|=d(X).
\]
The induction hypothesis applied to $Y$ supplies the required terminal
successor $Z$.  Finally apply the resulting assertion to the given initial
state $X_0$.
\end{proof}
